\documentclass[10pt,a4paper]{article}
\usepackage{amsmath,amssymb,amsthm,graphicx}
\usepackage{amsthm,amsfonts,amssymb,amsmath}
\newtheorem{theorem}{Theorem}[section]
\newtheorem{proposition}[theorem]{Proposition}
\newtheorem{lemma}[theorem]{Lemma}
\newtheorem{conjecture}[theorem]{Conjecture}
\newtheorem{fact}[theorem]{Fact}
\newtheorem{claim}[theorem]{Claim}
\newtheorem{corollary}[theorem]{Corollary}
\newtheorem{definition}[theorem]{Definition}
\newtheorem{problem}[theorem]{Problem}
\newtheorem{question}[theorem]{Question}
\newtheorem{observation}[theorem]{Observation}

\theoremstyle{definition}
\newtheorem{remark}[theorem]{Remark}
\newcommand{\su}{\subseteq}
\newcommand{\sm}{\setminus}
\newcommand{\Z}{\mathbb{Z}}
\newcommand{\HH}{\mathcal{H}}
\newcommand{\conv}{\operatorname{conv}}
\newcommand{\PP}{\mathcal{P}}
\title{An interval-wise partite matching family is bounded by the product of the pattern parameters}
\author{Lisa Sauermann, Dmitrii Zakharov}
\date{arxiv:2309.04813v4; Creative Commons Attribution 4.0 International}
\begin{document}
\maketitle
\noindent\textit{Editorial note.} The counted statement is the author's labelled lemma environment at source0.tex lines 359-362; its proof environment opens at line 365, two lines after \textbackslash{}end\{lemma\}, so the pairing is by adjacency of the statement with the following proof environment, and the retained proof range 365-377 is the whole environment.\par\medskip
\noindent\textit{Editorial note.} The preamble is built from the author's own package line, theorem-environment family and macro block (source lines 25-59); the journal document class of the source and its drawing packages are not retained because the retained material does not use them.\par\medskip
\noindent\textit{Editorial note.} Metadata provenance: the pinned source carries both a licence evidence file (source\_version 2309.04813v4, CC BY 4.0) and a source record file; the card's source\_version is the value recorded in the licence evidence. The source's own bibliography is inline in the manuscript and is not needed by the retained proof.\par\medskip
\noindent\textit{Editorial note.} The source carries no label-number dump and no counter values are imposed, so the printed numbering is the wrapper's default; the identity of the counted item is its author label obs-interval-wise.\par\medskip
\noindent\textit{Editorial note.} r002 note: the pre-publication validator requires either status=supported in the licence evidence or a retained-licence form with no status at all. This source's evidence file carries no acceptance record and none may be written on anyone's behalf, so r002 keeps no status and records a per-card licence reuse record (name, version\_url, prior\_license\_path, basis) pointing at the licence record already staged in the source directory form in a task-directory evidence copy that the assembler carries into the package: the source's staged record holds the official exact-version CC BY 4.0 observation. The original keys of the source evidence are retained unchanged, no new licence claim is made, and no acceptance record is written.\par\medskip
\noindent\textit{Editorial scope.} One original result of the pinned author manuscript, counted once: the product bound for interval-wise partite hypergraph matchings (author label obs-interval-wise).\par\medskip
\noindent\textit{Allowed standard prerequisites (editorial).}\begin{itemize}
\item Ordered hypergraph matching. A family of pairwise disjoint subsets of an ordered vertex set, with the induced order on each edge.. It is the family being counted.
\item Interval-wise partite family. There are cut points splitting the order into consecutive blocks so that every edge meets each block in exactly the prescribed number of vertices.. It is the structural hypothesis of the statement.
\item Pattern parameters. To each pattern of signs compatible with the ordered partition the construction attaches a parameter, and the product of those parameters is the bound.. They are the factors of the asserted estimate.
\end{itemize}\medskip
\begin{theorem} \label{thm-paramteters-upper-bound}
Let $r$ be a positive integer, and consider a positive integer $m_P$ for every collectable $r$-pattern $P$. Let $\HH$ be an ordered $r$-uniform hypergraph matching, and assume that for every collectable $r$-pattern $P$, every $P$-clique contained in $\HH$ has size at most $m_P$. Then the number of edges in $\HH$ is bounded by
\begin{equation}\label{eq_main_upper}
|\HH|\le 2^{r-1} \sum_{\lambda^{(1)} \succ \lambda^{(2)} \succ \dots \succ \lambda^{(r)}}\ \prod_{P\in \PP(\lambda^{(1)})\cup\dots\cup \PP(\lambda^{(r)})} m_{P},
\end{equation}
where the sum is over all sequences $\lambda^{(1)} \succ \lambda^{(2)} \succ \dots \succ \lambda^{(r)}$ of ordered partitions of $r$ (note that then for $s=1,\dots,r$, the ordered partition $\lambda^{(s)}$ must automatically have $s$ parts).
\end{theorem}

\label{sect-preliminaries}

\begin{theorem}[Multi-dimensional Erd\H{o}s--Szekeres Theorem]\label{thm_es}
Let $s$ be a positive integer. Furthermore, consider a positive integer $m_\tau$ for every function $\tau:\{1,\dots,s\}\to \{1,-1\}$ with $\tau(1)=1$. Let $Z\su \mathbb{R}^s$ be a collection of points, whose $|Z|\cdot s$ coordinates are all distinct. For each $\tau:\{1,\dots,s\}\to \{1,-1\}$ with $\tau(1)=1$, assume that every $\tau$-monotone sequence of points in $Z$ has length at most $m_\tau$. Then the number of points in $Z$ is bounded by
\[|Z|\le \prod_{\tau} m_{\tau},\]
where the product is over all functions $\tau:\{1,\dots,s\}\to \{1,-1\}$ with $\tau(1)=1$.

Furthermore, for any choice of positive integers $m_\tau$ for all functions $\tau:\{1,\dots,s\}\to \{1,-1\}$ with $\tau(1)=1$, there exists a collection of points $Z\su \mathbb{R}^s$ of size $|Z|=\prod_\tau m_\tau$ satisfying the conditions above.
\end{theorem}

\begin{lemma}\label{obs-interval-wise}
Let $r$ be a positive integer, and let $\lambda=(\lambda_1,\dots,\lambda_s)$ be an ordered partition of $r$. Furthermore, consider a positive integer $m_P$ for every $r$-pattern $P\in \PP(\lambda)$. Let $\HH$ be a finite collection of pairwise disjoint subsets $e\su \Z$, each of size $|e|=r$, such that $\HH$ is interval-wise $\lambda$-partite. For every $P\in \PP(\lambda)$, assume that every $P$-clique contained in $\HH$ has size at most $m_{P}$. Then we have
\[|\HH|\le \prod_{P\in \PP(\lambda)} m_{P}.\]
\end{lemma}

\begin{proof} 
    As $\HH$ is interval-wise $\lambda$-partite, there exist $x_0<x_1<\dots<x_{s}$ such that for each $i=1,\dots,s$, we have $|e\cap (x_{i-1},x_i)|=\lambda_i>0$ for every $e\in \HH$ and the intervals $\conv(e\cap (x_{i-1},x_i))$ are pairwise disjoint for all $e\in \HH$. Now, for every $e\in \HH$, let us define a point $z(e)\in \mathbb{R}^s$ by taking
    \[z(e)=(\min(e\cap (x_{0},x_1)), \min(e\cap (x_{1},x_2)), \dots, \min(e\cap (x_{s-1},x_s))),\]
    and define $Z=\{ z(e)\mid e\in \HH\}\su \mathbb{R}^s$. As the sets $e\in \HH$ are pairwise disjoint, the all the coordinates of the points $z(e)$ for $e\in \HH$ are distinct.
    
    For every function $\tau:\{1,\dots,s\}\to \{1,-1\}$ with $\tau(1)=1$, let us define $m_\tau=m_{P(\lambda,\tau)}$ (recall that $P(\lambda,\tau)\in \PP(\lambda)$ was defined at the end of Section~\ref{sect-preliminaries}). We claim that then every  $\tau$-monotone sequence of points in $Z\su \mathbb{R}^s$ has length at most $m_\tau$.

    Indeed, for any $\rho:\{1,\dots,s\}\to \{1,-1\}$ with $\tau(1)=1$, consider $e^{(1)},e^{(2)},\dots,e^{(t)}\in \HH$ such that the point sequence $z(e^{(1)}),z(e^{(2)}),\dots,z(e^{(t)})$ is a $\tau$-monotone sequence in $Z \su \mathbb{R}^s$. Then for any index $i\in \{1,\dots,s\}$ with $\tau(i)=1$ we have  $\min(e^{(1)}\cap (x_{i-1},x_i))< \dots <\min(e^{(t)}\cap (x_{i-1},x_i))$, and in particular $\min(e^{(\ell)}\cap (x_{i-1},x_i))<\min(e^{(\ell')}\cap (x_{i-1},x_i))$ for any $1\le \ell<\ell'\le t$. Since furthermore the intervals $\conv(e^{(\ell)}\cap (x_{i-1},x_i))$ and $\conv(e^{(\ell')}\cap (x_{i-1},x_i))$ are disjoint, this means that all elements of $e^{(\ell)}\cap (x_{i-1},x_i)$ are smaller than all elements of $e^{(\ell')}\cap (x_{i-1},x_i)$. Similarly, for any $i\in \{1,\dots,s\}$ with $\tau(i)=-1$ and any $1\le \ell<\ell'\le t$, all elements of $e^{(\ell)}\cap (x_{i-1},x_i)$ are larger than all elements of $e^{(\ell')}\cap (x_{i-1},x_i)$. Since $x_0<x_1<\dots<x_{s}$ and $|e^{(\ell)}\cap (x_{i-1},x_i)|=|e^{(\ell')}\cap (x_{i-1},x_i)|=\lambda_i$ for all  $i=1,\dots,s$, this shows that  $e^{(\ell)}$ and $e^{(\ell')}$ form the pattern $P(\lambda,\tau)$ for all $1\le \ell<\ell'\le t$. Hence $e^{(1)},e^{(2)},\dots,e^{(t)}\in \HH$ form a $P(\lambda,\tau)$-clique and therefore $t\le m_{P(\lambda,\tau)}=m_\rho$. So every $\rho$-monotone sequence in $Z$ indeed has length at most $m_\rho$.

    Applying Theorem~\ref{thm_es} now gives
    \[|\HH|=|Z|\le \prod_{\substack{\tau:\{1,\dots,s\}\to \{1,-1\}\\ \tau(1)=1}} m_{\tau}=\prod_{\substack{\tau:\{1,\dots,s\}\to \{1,-1\}\\ \tau(1)=1}} m_{P(\lambda,\tau)}=\prod_{P\in \PP(\lambda)} m_{P},\]
    as desired.
\end{proof}

\end{document}
